% Shared preamble for the in-class practice handouts and answer keys.
% Each main file sets \chapterlabel and \handoutkind; compile with xelatex.
\documentclass[UTF8, scheme=plain, fontset=fandol, 11pt, a4paper]{ctexart}

% Same CJK fonts as elegantbook-zh-tw.cls when installed (e.g. on Overleaf);
% otherwise fall back to the macOS Traditional Chinese fonts (Fandol lacks many TC glyphs).
\IfFontExistsTF{AR PL UMing TW MBE}{%
    \setCJKmainfont[ItalicFont=AR PL UKai TW, BoldFont=cwTeXHeiBold]{AR PL UMing TW MBE}%
}{\IfFontExistsTF{Songti TC}{%
    \setCJKmainfont[ItalicFont=Kaiti TC, BoldFont=Heiti TC]{Songti TC}%
}{}}

\usepackage[margin=2cm, top=2.2cm, bottom=2cm]{geometry}
\usepackage{amsmath, amssymb}
\usepackage{xcolor}
\usepackage{enumitem}
\usepackage{fancyhdr}
\usepackage[most]{tcolorbox}

\definecolor{structurecolor}{RGB}{60,113,183}
\definecolor{third}{RGB}{0,160,152}

\newcommand{\PP}{\mathbb{P}}
\newcommand{\EE}{\mathbb{E}}
\newcommand{\NN}{\mathcal{N}}
\newcommand{\ZZ}{\mathcal{Z}}
\newcommand{\var}{var}
\newcommand{\sd}{sd}

\setlength{\parindent}{0pt}
\renewcommand{\baselinestretch}{1.25}

\pagestyle{fancy}
\fancyhf{}
\renewcommand{\headrule}{\color{structurecolor}\hrule width\textwidth}
\newcommand{\chapterlabel}{}
\fancyhead[L]{\small 生物統計\quad \chapterlabel}
\fancyhead[R]{\small\handoutkind}
\fancyfoot[C]{\small\thepage}

% \grouptitle{組別}{難度星數}{主題}
\newcommand{\grouptitle}[3]{%
    {\Large\bfseries\color{structurecolor} 第 #1 組\quad #3}\hfill{\color{third}#2}\par\medskip}

% Scenario box
\newtcolorbox{scenario}{enhanced, breakable, colback=structurecolor!5, colframe=structurecolor,
    boxrule=0.6pt, arc=2pt, left=6pt, right=6pt, top=4pt, bottom=4pt,
    fonttitle=\bfseries, title=情境}

\newcommand{\stars}[1]{\small\count255=0 \loop\ifnum\count255<#1 $\bigstar$\advance\count255 by 1 \repeat}

% \scenariotitle{編號}{主題}{章節}: heading for a multi-part scenario
\newcommand{\scenariotitle}[3]{%
    {\Large\bfseries\color{structurecolor} 情境 #1\quad #2}\hfill{\color{third}\small #3}\par\medskip}

% Numbered subproblems with a difficulty marker: \sub{星數} text
\newenvironment{subproblems}{\begin{enumerate}[leftmargin=*, label=\textbf{\arabic*.}, itemsep=8pt]}{\end{enumerate}}
\newcommand{\sub}[1]{\item {\color{third}\stars{#1}}\quad}

\renewcommand{\chapterlabel}{第 4 章\ 隨機變數與機率分布}
\newcommand{\handoutkind}{課堂分組練習\ 參考解答}

\newenvironment{answers}{\begin{enumerate}[leftmargin=*, label=\textbf{\arabic*.}, itemsep=6pt]}{\end{enumerate}}

\begin{document}

%% ---------------------------------------------------------------- 第 1 組
\grouptitle{1}{\stars{1}}{疫苗注射部位痠痛}
\begin{answers}
    \item $X \sim Binomial(5, 0.2)$。假設：各人是否痠痛互相獨立，且每人痠痛機率相同。
    \item $\PP(X=1) = \binom{5}{1}(0.2)(0.8)^4 = 0.4096$；\quad $\PP(X\ge 1) = 1 - \PP(X=0) = 1-0.8^5 = 0.6723$。
    \item $\EE(X) = np = 1$，$\var(X) = np(1-p) = 0.8$。\\
          $\EE(Y) = 50\,\EE(X) + 200 = 250$ 元；$\var(Y) = 50^2 \times 0.8 = 2000$，$\sd(Y) = 50\sqrt{0.8} \approx 44.7$ 元。固定成本只平移分布，不影響分散程度。
    \item $z = (38.0-37.0)/0.5 = 2$，$\PP(\ZZ \ge 2) = 0.0228$。
    \item $z_{0.025} = 1.96$，門檻為 $37.0 + 1.96 \times 0.5 = 37.98$°C。
    \item 不合理。同一家人有共同的體質與環境，也可能互相暗示，\textbf{獨立性}最可能被違反；「每人機率相同」也未必成立。此時痠痛人數會比二項式分布更容易出現「全有或全無」，變異數較大。
\end{answers}
\newpage

%% ---------------------------------------------------------------- 第 2 組
\grouptitle{2}{\stars{2}}{高血壓服藥遵從}
\begin{answers}
    \item $X \sim Binomial(10, 0.9)$，$p>0.5$ 故為左偏；$W \sim Binomial(10, 0.1)$，$p<0.5$ 故為右偏。兩者的機率質量函數互為鏡像。
    \item $\PP(X \ge 9) = \binom{10}{9}(0.9)^9(0.1) + (0.9)^{10} = 0.3874 + 0.3487 = 0.7361$。\\
          $X \ge 9$ 與 $W \le 1$ 是同一個事件，所以 $\PP(W \le 1)$ 的算式與答案完全相同。
    \item $\EE(X) = 9$，$\var(X) = 10(0.9)(0.1) = 0.9$；$\EE(W) = 1$，$\var(W) = 10(0.1)(0.9) = 0.9$。\\
          $W = -X + 10$，故 $\var(W) = (-1)^2\var(X) = \var(X)$，$\EE(W) = -\EE(X) + 10$。
    \item $z_1 = (122-128)/12 = -0.5$，$z_2 = (140-128)/12 = 1$。\\
          $\PP(-0.5 \le \ZZ \le 1) = 0.8413 - 0.3085 = 0.5328$。
    \item $z_{0.05} = 1.645$，門檻為 $128 - 1.645 \times 12 = 108.3$ mmHg。
    \item $\var(X \pm Y) = \var(X) + \var(Y)$ 的前提是 $X$ 與 $Y$ \textbf{獨立}。這裡 $W$ 完全由 $X$ 決定，兩者完全相依，公式不適用。期望值的相加則不需要獨立：$\EE(X+W) = 9 + 1 = 10$ 仍然成立。
\end{answers}
\newpage

%% ---------------------------------------------------------------- 第 3 組
\grouptitle{3}{\stars{3}}{急診胸痛病人}
\begin{answers}
    \item 每小時 3 位，20 分鐘為 $1/3$ 小時，故 $\lambda = 1$，$X \sim Poisson(1)$。條件：(1) 事件之間獨立；(2) 發生頻率不隨時間改變；(3) 同一時間點不會發生兩次以上事件。
    \item $\PP(X=0) = e^{-1} = 0.3679$；\quad $\PP(X \ge 2) = 1 - \PP(X=0) - \PP(X=1) = 1 - 2e^{-1} = 0.2642$。
    \item 8 小時：$\lambda = 24$，期望值 $24$，標準差 $\sqrt{24} \approx 4.90$。\\
          16 小時：$\lambda = 48$，標準差 $\sqrt{48} \approx 6.93$，只變為 $\sqrt{2}$ 倍。期望值與\textbf{變異數}隨時間等比例增加，標準差則否。
    \item $z = (10-8)/2 = 1$，$\PP(\ZZ \le 1) = 0.8413$。
    \item 需要 $\PP(X \le 10) = 0.975$，即 $(10-\mu)/2 = 1.96$，$\mu = 10 - 1.96 \times 2 = 6.08$ 分鐘。
    \item 條件 (2) 最容易不成立（白天與深夜、平日與連假的速率不同）；條件 (1)(3) 在群聚事件（如同一事故多人送醫）時不成立。\\
          若每一天的速率固定但各天不同，把許多天的資料混在一起就是「來自多個速率參數不同的母體」，變異數會\textbf{大於}平均值（過度離散）。注意：若只是同一天內速率隨時段變動、但每天的變動型態相同，則每日總人數仍為卜瓦松分布。
\end{answers}
\newpage

%% ---------------------------------------------------------------- 第 4 組
\grouptitle{4}{\stars{4}}{血紅素與貧血}
\begin{answers}
    \item $B \sim Bernoulli(p)$，$X \sim Binomial(8, p)$，其中 $p$ 是常態分布在 12 g/dL 以下的曲線下面積（第 4 題）。連續變數經過二分後成為白努利變數，再加總即為二項式變數。
    \item $\PP(X \ge 2) = 1 - (1-p)^8 - 8p(1-p)^7 = 1 - 0.4095 - 0.3869 = 0.2039$（以 $p = 0.1056$ 計算）。
    \item $\EE(X) = 8p = 0.845$，$\var(X) = 8p(1-p) = 0.756$。\\
          抽 80 人時，期望值變為 10 倍，變異數變為 10 倍，故標準差變為 $\sqrt{10} \approx 3.16$ 倍。
    \item $z = (12-13.5)/1.2 = -1.25$，$p = \PP(\ZZ \le -1.25) = 0.1056$。
    \item 需要 $\PP(X \ge x) = 0.80$，$z_{0.80} = -z_{0.20} = -0.8416$，門檻為 $13.5 - 0.8416 \times 1.2 = 12.49$ g/dL。
    \item 損失了\textbf{嚴重程度}的資訊：血紅素 11.9 與 8.0 都被編碼為 $1$。\\
          平均值降為 12.9 時，$z = (12-12.9)/1.2 = -0.75$，貧血比例為 $0.2266$，是原本的兩倍以上。平均值只移動半個標準差，尾端比例卻倍增，變化並非線性。
\end{answers}
\newpage

%% ---------------------------------------------------------------- 第 5 組
\grouptitle{5}{\stars{5}}{癲癇發作}
\begin{answers}
    \item 不是卜瓦松分布，而是兩個卜瓦松分布的混合：
          \[f_X(x) = 0.8 \cdot \frac{0.5^x e^{-0.5}}{x!} + 0.2 \cdot \frac{4^x e^{-4}}{x!}, \qquad x \in \{0,1,2,\dots\}\]
    \item $\PP(X=0) = 0.8e^{-0.5} + 0.2e^{-4} = 0.4852 + 0.0037 = 0.4889$。\\
          $Poisson(1.2)$ 的 $\PP(X=0) = e^{-1.2} = 0.3012$。混合母體中零的比例明顯偏多。
    \item $\EE(X) = 0.8 \times 0.5 + 0.2 \times 4 = 1.2$。\\
          各子群 $\EE(X^2) = \lambda + \lambda^2$：控制良好組為 $0.75$，控制不佳組為 $20$。\\
          $\EE(X^2) = 0.8 \times 0.75 + 0.2 \times 20 = 4.6$，$\var(X) = 4.6 - 1.2^2 = 3.16 > 1.2$。變異數大於期望值。
    \item $z_1 = (50-70)/15 = -1.33$，$z_2 = (100-70)/15 = 2$。\\
          $\PP(-1.33 \le \ZZ \le 2) = 0.9772 - 0.0918 = 0.885$（不四捨五入 $z_1$ 則為 $0.886$）。
    \item $z_{0.10} = 1.2816$，$z_{0.05} = 1.645$，故
          \[\mu - 1.2816\,\sigma = 50, \qquad \mu + 1.645\,\sigma = 100\]
          相減得 $\sigma = 50/2.9266 = 17.1$，代回得 $\mu = 71.9$ $\mu$g/mL。
    \item 樣本變異數遠大於平均值，表示資料不符合單一卜瓦松分布，母體很可能由速率不同的子群組成（過度離散）。\\
          由貝氏定理：$\PP(\text{控制不佳} \mid X=0) = \dfrac{0.2e^{-4}}{0.4889} = \dfrac{0.0037}{0.4889} \approx 0.0075$。
\end{answers}
\newpage

%% ---------------------------------------------------------------- 第 6 組
\grouptitle{6}{\stars{6}}{降壓藥臨床試驗}
\begin{answers}
    \item 新藥組反應率 $p_T = \PP(D_T \ge 10) = \PP(\ZZ \ge -0.25) = 0.5987$；\\
          安慰劑組反應率 $p_C = \PP(D_C \ge 10) = \PP(\ZZ \ge 0.625) = 0.2660$。\\
          $X_T \sim Binomial(6, 0.5987)$，$X_C \sim Binomial(6, 0.2660)$，兩者獨立。
    \item 兩組獨立，機率相乘：$\PP(X_T = 6)\,\PP(X_C = 0) = 0.5987^6 \times 0.7340^6 = 0.0461 \times 0.1564 = 0.0072$。
    \item $\EE(X_T - X_C) = 6(0.5987 - 0.2660) = 2.00$。\\
          $\var(X_T - X_C) = \var(X_T) + \var(X_C) = 6(0.5987)(0.4013) + 6(0.2660)(0.7340) = 1.442 + 1.171 = 2.61$。\\
          相減時期望值相減，但變異數仍是\textbf{相加}。
    \item $D_T - D_C \sim \NN(12-5,\; 8^2+8^2) = \NN(7, 128)$，標準差 $\sqrt{128} = 11.31$。\\
          $\PP(D_T - D_C > 0) = \PP(\ZZ > -7/11.31) = \PP(\ZZ > -0.62) = 0.732$。
    \item 門檻為 $5 + 1.645 \times 8 = 18.16$ mmHg。\\
          新藥組：$z = (18.16-12)/8 = 0.77$，反應率 $\PP(\ZZ \ge 0.77) = 0.221$。
    \item 不成立。同一位病人的兩次測量通常為正相關，不再獨立。正相關時，個人特質造成的變異在相減時互相抵消，差值的變異數會\textbf{更小}。因此在相同人數下，治療效果（平均差 7 mmHg）相對於雜訊更明顯，更容易偵測到差異——這就是配對設計的優勢（第 8 章）。
\end{answers}

\end{document}
